% !TeX root = ../初等数学笔记.tex
% 由主文档通过 \input 载入；请编译主文档。
\setcounter{chapter}{11} % 保留原书章号。
\chapter{有限问题中的证明方法}
\label{chap:em-discrete}
有限问题常常不需要复杂公式，而需要找到对象之间的组织方式。抽屉原理保证某种重复，不变量排除不可达状态，极端原理抓住最特殊对象，双重计数则把局部关系汇总为整体等式。

\section{抽屉原理与存在性}
\begin{theorem}[抽屉原理]
将 $N$ 个对象放入 $m\ge1$ 个盒子，至少一个盒子含有 $\lceil N/m\rceil$ 个或更多对象；$\lceil t\rceil$ 表示不小于 $t$ 的最小整数。
\end{theorem}
\begin{proof}
若每盒至多 $\lceil N/m\rceil-1$ 个，则总数至多为 $m(\lceil N/m\rceil-1)<N$，与对象总数矛盾。
\end{proof}
任取 $m+1$ 个整数，以模 $m$ 的余数作为盒子，则其中两数之差被 $m$ 整除。关键不是逐一计算差，而是识别只有 $m$ 种余数。存在性结论未必告诉我们具体是哪一对，若要找出对象，还需实际分组或搜索。

\section{不变量与单调量}
\emph{不变量}（invariant）是每次允许操作都保持的量。若初始与目标状态的不变量不同，目标就不可达。把棋盘按黑白交替染色，一块覆盖相邻两格的骨牌必覆盖一黑一白；从 $8\times8$ 棋盘删去同色的两个对角格后，黑白格数量相差 $2$，因此不能用骨牌无重叠覆盖。这一证明不必枚举所有铺法。

严格变化的量也能证明终止，但还需取值范围不能无限下降。欧几里得算法的余数是严格下降的非负整数，所以有限步终止；正实数 $1,1/2,1/4,\ldots$ 却能永远下降，故“每步减小”本身不足够。不变量一般只给必要条件：颜色数量相等不能单独保证任意形状都能铺满。

\section{极端原理与最小反例}
有限非空集合的实数值必能取得最大和最小值；正整数集合即使无限，只要非空也有最小元素。这使我们可以选择最小反例，再从它构造更小反例以得到矛盾，称为无限递降法。

\begin{example}
证明不存在正整数 $a,b,c$ 使 $a^2+b^2=3c^2$。
\end{example}
\begin{solve}
若存在，取 $c$ 最小的一组解。平方模 $3$ 只可能为 $0$ 或 $1$，两平方和模 $3$ 为零迫使 $a,b$ 都被 $3$ 整除，写成 $a=3u,b=3v$。代入得 $c^2=3(u^2+v^2)$，所以 $c=3w$。再约去 $9$ 得 $u^2+v^2=3w^2$，其中 $0<w=c/3<c$，与最小性矛盾。
\end{solve}
选择极端量时必须说明它存在，且构造后的对象仍满足全部条件；缺少后一检查，递降论证就不完整。

\section{双重计数与简单图论选读}
用两种方式数同一集合称为\emph{双重计数}（double counting）。有限无向简单图由顶点集与连接不同顶点的无向边组成，不含自环与重边；顶点的度是与它相接的边数。每条边恰有两个端点，所以
\[
 \sum_{v}\deg(v)=2|E|,
\]
其中 $E$ 为边集。左端奇偶性由奇度顶点数量决定，而右端为偶数，故奇度顶点个数必为偶数。

同一思想还能证明 $\binom n2=n(n-1)/2$：先数有序不同元素对得 $n(n-1)$，每个无序对被计两次。图的画法可以变化，边是否交叉及画得多长通常不属于图的组合数据；解释模型时应避免给图形增加未声明的几何条件。

\section{递推计数、构造与算法描述}
设 $T_n$ 为用 $1\times2$ 骨牌铺满 $2\times n$ 长方形的方法数。空区域有一种铺法，故 $T_0=1$，且 $T_1=1$。考察最左侧，或者放一块竖骨牌留下 $2\times(n-1)$ 区域，或者放两块横骨牌留下 $2\times(n-2)$ 区域，两类互斥且穷尽，因此
\[
 T_n=T_{n-1}+T_{n-2},\qquad n\ge2.
\]
依次维护相邻两个数即可算出 $T_n$，需要 $O(n)$ 次整数加法和常数个整数存储位置；若计入整数位数增长，则总位运算成本不是简单的 $O(n)$。

构造题除了证明目标不被不变量排除，还要给出合法操作序列，并说明何时结束。算法的正确性可以按循环步骤归纳，复杂度则需说明数的是运算次数、比较次数还是实际位操作。

\section{例题、习题与学习检查}
\begin{example}
单位正方形内任取五个点，证明其中两点的距离不超过 $\sqrt2/2$。
\end{example}
\begin{solve}
将正方形分成边长 $1/2$ 的四个小正方形，边界点按照固定规则归入其中一个。由抽屉原理有两点归入同一小方形，它们的距离不超过该方形对角线 $\sqrt{(1/2)^2+(1/2)^2}=\sqrt2/2$。
\end{solve}
\begin{enumerate}
\item 任取 $7$ 个整数，证明其中两数之差为 $6$ 的倍数。提示：按模 $6$ 余数分组。
\item 某图的顶点度数为 $1,1,2,2,4$，若该图存在，边数是多少？答：$5$；握手关系只给必要计数约束。
\item 计算 $T_5$ 并说明不能把两个第一步分支相乘。答：$T_5=8$，它们是互斥选择而非顺次步骤。
\end{enumerate}

